\documentclass[11pt,reqno]{amsart}
\usepackage{amsmath,amssymb,graphicx,amscd}
\usepackage{mathrsfs}
\usepackage{enumerate}
% pour les S des spheres
\usepackage{mathrsfs}
\usepackage{dsfont}
\usepackage{stmaryrd}
\usepackage{wrapfig}
\usepackage{pstricks}

% ----------------------------------------------------------------
\vfuzz2pt % Don't report over-full v-boxes if over-edge is small
\hfuzz2pt % Don't report over-full h-boxes if over-edge is small

\newcommand{\prob}{{\mathbb{P}}}

\renewcommand{\S}{\mathscr{S}}  % complex unit sphere

\newcommand{\CC}{{\mathbb C}}
\newcommand{\HH}{{\mathbb H}}
\newcommand{\RR}{{\mathbb R}}
\newcommand{\QQ}{{\mathbb{Q}}}
\newcommand{\NN}{{\mathbb{N}}}
\newcommand{\ZZ}{{\mathbb{Z}}}
\renewcommand{\P}{{\mathbb{P}}}
\renewcommand{\span}{{\mathrm{span}}}
\newcommand{\sgn}{\operatorname{{sgn}}}
 \DeclareMathOperator{\Prob}{\mathbb{P}}   %probability
 \DeclareMathOperator{\E}{\mathbb{E}}      %expectation
 \DeclareMathOperator{\cov}{cov}             %logarithmic integral

 \newcommand{\I}{1\!\!1}                   %Indicator function


 \newcommand{\dd}{{\mathrm{d}}}            %differential
 \newcommand{\ii}{{\mathrm{i}}}
  \newcommand{\Id}{{\mathrm{Id}}}
\newcommand{\law}{\overset{\mbox{\rm \scriptsize law}}{=}}
\newcommand{\convlaw}{\overset{\mbox{\rm \scriptsize law}}{\longrightarrow}}


\renewcommand{\Re}{{\mathfrak{Re}}}
\renewcommand{\Im}{{\mathfrak{Im}}}

% THEOREMS -------------------------------------------------------

\newtheorem{thm}{Theorem}[section]
\newtheorem{cor}[thm]{Corollary}
\newtheorem{lem}[thm]{Lemma}
\newtheorem{prop}[thm]{Proposition}
\theoremstyle{definition}
\newtheorem{defn}[thm]{Definition}
\theoremstyle{remark}
\newtheorem*{rem}{Remark}
 \newcommand{\rd}{{\mathrm{d}}}            %differential
\newcommand{\1}{\mathds{1}}


\numberwithin{equation}{section}



\begin{document}

~\vspace{-3cm}
\begin{center}
{\bf Probability, homework 3 due September 23.}
\end{center}
\vspace{1cm}
\small



\noindent{\bf Exercise 1}.
Let $(\Omega,\mathcal{A},\Prob)$ be a probability space. Prove that if $A\cap B=\emptyset$ and $A,B$ are independent, then $\Prob(A)=0$ or $\Prob(B)=0$.\\



\noindent{\bf Exercise 2}. Let $X$ be a nonnegative random variable with null expectation. Prove that it is $0$ almost surely.\\



\noindent{\bf Exercise 3}. Calculate $\E(X)$ for the following probability measures $\mathbb{P}^X$.
\begin{enumerate}[(i)]
\item $\mathbb{P}^X$ has Gaussian density $\frac{1}{\sqrt{2\pi}\sigma}e^{-(x-\mu)^2/(2\sigma^2)}$, for some $\sigma>0$ and $\mu\in\RR$;
\item $\mathbb{P}^X$ has exponential dentity $\lambda e^{-\lambda x}\mathds{1}_{x\geq 0}$ for some $\lambda>0$;
\item $\mathbb{P}^X=p \delta_a+q\delta_b$ where $p+q=1$, $p,q\geq 0$ and $a,b\in\RR$;
\item $\mathbb{P}^X$ is the Poisson distribution: $\mathbb{P}^X(\{n\})=e^{-\lambda}\frac{\lambda^n}{n!}$ for any integer $n\geq 0$, for some $\lambda>0$.\\
\end{enumerate}


\noindent{\bf Exercise 4}.
Let $X$ be a standard Gaussian random variable. What is the density of $1/X^2$?\\



\noindent{\bf Exercise 5}.
Let $X$ be uniformly distributed on $[0,1]$ and $\lambda>0$. Show that $-\lambda^{-1}\log X$ has the same distribution as an exponential random variable with parameter $\lambda$.\\



\noindent{\bf Exercise 6}. A samoura\"i wants to create a triangle with a (rigid) spaghetti. With his saber, he cuts this spaghetti on two places, chosen uniformly and independently along this traditional pasta. What is the probability that he can create a triangle with sides these three pieces of spaghetti?\\

\noindent{\bf Exercise 7}. Assume that $X_1,X_2,\dots$ are independent random variables uniformly distributed on $[0,1]$.
Let $Y^{(n)}=n \inf\{X_i,1\leq i\leq n\}$. Prove that it converges weakly to an exponential random variable, i.e.
for any continuous bounded function $f:\RR^+\to\RR$,
$$
\E\left(f(Y^{(n)})\right)\underset{n\to\infty}{\longrightarrow}\int_{\RR^+}f(u)e^{-u}{\rm d} u.
$$

\noindent{\bf Exercise 8}.  Let $f:[0,1]\to\mathbb{R}$ be continuous.  For any $y\in\mathbb{R}$, let $N(y)\in\mathbb{R}\cup\{\infty\}$ be the number of solutions to $f(x)=y$. Prove that $N$ is measurable. \\

\noindent{\bf Exercise 9}.  Let $f:\mathbb{R}\to\mathbb{R}$ be additive ($f(a+b)=f(a)+f(b)$ for all $a$, $b$) and measurable.  Prove that it is  linear.

What if $f$ is not assumed to be measurable? What about $f:\mathbb{R}^d\to\mathbb{R}$?

\end{document} 